\section{A Unique Solution}
It is essential to verify if the equation system given in \cref{eqn:Compact_Eq_Sys} has a
\emph{unique} solution. From the linear algebra we know that such a system has always a
unique solution if and only if determinant of the coefficient matrix is nonzero.

\begin{definition} \label{def:Positive_Definite}
	A real square matrix $\gls{matA} \in \gls{Rnn}$ is said to be \gls{posDef} if and only it is symmetrical and
	$\gls{v}^{T} \gls{matA} \gls{v} > 0$, for an arbitrary vector $\gls{v} \in \mathbb{R}^{n \times 1}$
	\begin{subequations}
		\begin{align}
			\gls{matA} &= \gls{matAT}, \label{eqn:Pos_Def_Prop_1} \\
			\gls{v}^{T} \gls{matA} \gls{v} &= \sum \limits_{i = 1}^{n} \sum \limits_{j = 1}^{n} v_{i}  \gls{matAij} v_{j} > 0,
				\quad \forall \mathbf{v} \in \mathbb{R}^{n \times 1} \label{eqn:Pos_Def_Prop_2}
		\end{align}
	\end{subequations}
\end{definition}
\begin{lemma} \label{lem:Pos_Def_Non_Singular}
	Suppose $\gls{matA} \in \gls{Rnn}$ is positive-definite. Then
	\begin{equation}
		\gls{detA} \neq 0.
	\end{equation}
\end{lemma}
For a proof of \cref{lem:Pos_Def_Non_Singular} any textbook in linear algebra may be conferred. In the following we will show
that a \gls{posDef} matrix is \gls{nonSingular} for the simple case of $\mathbf{A} \in \mathbb{R}^{2 \times 2}$. \\
\begin{example}
Suppose that $\gls{matA} \in \glsentryname{R22}$ is a positive-definite matrix given by
\begin{equation}
	\gls{matA} =
	\begin{bmatrix}
		a & b \\
		b & c	
	\end{bmatrix}.
\end{equation}
Then note that
\begin{equation}
\gls{detA} =
	\begin{vmatrix}
		a & b \\
		b & c	
	\end{vmatrix} = a c - b^{2}.
\end{equation}
Now define \gls{v} $\in$ \glsentryname{R21} by
\begin{equation}
	\gls{v} = \begin{bmatrix}	x \\ y	\end{bmatrix},
\end{equation}
and calculate
\begin{equation}	\label{eqn:Temp_Pos_Def_Exp}
	\begin{aligned}
		\gls{v}^T \gls{matA} \gls{v} &=
			\begin{bmatrix} x & y \end{bmatrix}
			\begin{bmatrix} a & b \\ b & c \end{bmatrix}
			\begin{bmatrix} x \\ y \end{bmatrix} \\
		&=	\begin{bmatrix} ax + by & bx + cy \end{bmatrix}
			\begin{bmatrix} x \\ y \end{bmatrix} \\
		&=	ax^{2} + 2bxy + cy^{2}.
	\end{aligned}
\end{equation}
Factoring and rearranging the terms in \cref{eqn:Temp_Pos_Def_Exp} we get
\begin{equation}	\label{eqn:Final_Pos_Def_Exp}
	\begin{aligned}
		\gls{v}^T \gls{matA} \gls{v} &= a \left( x^{2} + \frac{2b}{a} xy + \frac{c}{a} y^{2} \right) \\
		&= a \left( x + \frac{b}{a}y \right)^2 + \frac{c}{a} y^2 - \frac{b^2}{a^2} y^2 \\
		&= a \left( x + \frac{b}{a}y \right)^2 + \frac{ac - b^2}{a^2} y^2 > 0.
	\end{aligned}
\end{equation}
Note that the last inequality is a consequence of the positive-definite property of $\mathbf{A}$. Sufficient
condition to satisfy \cref{eqn:Final_Pos_Def_Exp} is
\begin{equation}
	a > 0, \quad ac - b^2 > 0 \Rightarrow \det(\mathbf{A}) > 0.
\end{equation}
A \gls{posDef} matrix $\mathbf{A}$ $\in \mathbb{R}^{2 \times 2}$ is then non-singular.
\end{example}
\begin{remark}
	A square matrix \gls{matA} $\in$ \gls{Rnn} may be \gls{nonSingular}, $\gls{detA} \neq 0$ without being \gls{posDef}, but
	if the matrix is \gls{posDef} then it is always \gls{nonSingular}, $\gls{detA} \neq 0$. On the other hand, if a matrix
	$\gls{matA}$ is \gls{singular}, it cannot be \gls{posDef}. These statements may be summarized as
	\begin{subequations}
		\begin{align} 
			\gls{matA} \text{ is positive-definite} & \Longrightarrow \gls{detA} \neq 0, \\
			\gls{detA} = 0 & \Longrightarrow \gls{matA} \text{ is not positive-definite}.
		\end{align}
		Thus positive definiteness is a more restrict condition than non-singularity.
	\end{subequations}
\end{remark}

Before imposing the essential boundary condition, the assembled stiffness matrix is only positive semidefinite: the constant function belongs to its nullspace. Uniqueness follows after restricting the discrete space to functions satisfying the homogeneous essential condition at the left boundary.

Let $\mathbf{K}_0$ denote the stiffness matrix on the free degrees of freedom. According to \cref{eqn:Stiffness_Matrix},
\begin{equation}
	\mathbf{K}_0 = \mathbf{K}_0^T.
\end{equation}
For any nonzero free coefficient vector $\mathbf{v}$, let $w_h$ be its associated discrete function. Then $w_h(0)=0$ and
\begin{equation}
	\mathbf{v}^{T}\mathbf{K}_0\mathbf{v}
	= \int_0^1 q(x)\,[w_h'(x)]^2 \dd x > 0. \label{eqn:Pos_Def_Simplified}
\end{equation}
Indeed, equality could hold only if $w_h$ were constant; together with $w_h(0)=0$, this would imply $w_h=0$ and hence $\mathbf{v}=0$. Thus $\mathbf{K}_0$ is positive definite and the constrained Galerkin solution is unique.

The implementation enforces the left value by replacing one matrix row. That full algebraic matrix is nonsymmetric, but it is equivalent to solving the symmetric positive-definite reduced system for the free degrees of freedom.
%
